\section*{Chapitre \ref{ch:b1:transient-regimes} --- Régimes transitoires~: premier et second ordre}

\begin{solution}{exo:b1:transient-regimes:1}
$\tau = RC = 10^4 \times 47\times10^{-9} = \qty{0.47}{ms}$~; $u_C(\tau)/E
= 1 - \eu^{-1} = 0.63$~; pour $99\%$~: $\eu^{-t/\tau} = 0.01$, $t = \tau\ln 100
= 4.6\tau = \qty{2.2}{ms}$.
\end{solution}

\begin{solution}{exo:b1:transient-regimes:2}
$\tau = L/R = \qty{4.0}{ms}$~; $i_\infty = E/R = \qty{2.0}{A}$~; $i(\tau) =
0.63 \times 2.0 = \qty{1.26}{A}$~; en $0^+$, $i = 0$ (continuité), donc
$u_L = E - Ri = \qty{10}{V}$.
\end{solution}

\begin{solution}{exo:b1:transient-regimes:3}
$\omega_0 = 1/\sqrt{LC} = \qty{3.16e4}{rad/s}$, $f_0 = \qty{5.03}{kHz}$~;
$Q = \sqrt{L/C}/R = 316/100 = 3.2 > \tfrac12$~: pseudo-périodique~;
$\tau = 2Q/\omega_0 = \qty{0.20}{ms}$~; $\omega = \omega_0\sqrt{1 - 1/4Q^2}
= 0.987\omega_0$, $T = \qty{0.20}{ms}$.
\end{solution}

\begin{solution}{exo:b1:transient-regimes:4}
$R_c = 2\sqrt{L/C} = \qty{632}{\ohm}$, $Q = \tfrac12$. À \qty{2}{k\ohm},
$Q = 0.16$~: apériodique, une reptation lente gouvernée par la petite
racine $\abs{r_+} \approx \omega_0 Q = \qty{5e3}{s^{-1}}$ (échelle de
temps \qty{0.2}{ms}, mais sans oscillation).
\end{solution}

\begin{solution}{exo:b1:transient-regimes:5}
$u = U_0\eu^{-t/\tau}$ est divisée par deux lorsque $\eu^{-t/\tau} = \tfrac12$~: $t_{1/2}
= \tau\ln 2$. $\tau = 3.5/0.693 = \qty{5.05}{ms}$, $R = \tau/C =
\qty{5.1}{k\ohm}$.
\end{solution}

\begin{solution}{exo:b1:transient-regimes:6}
$i = (E/R)\eu^{-t/\tau}$~: $\int_0^\infty Ri^2\,\dd t = (E^2/R)(\tau/2) =
\tfrac12 CE^2$, sans trace de $R$. La source fournit $CE^2$~: rendement
de $50\%$ quelle que soit la résistance. Mieux~: charger à travers une
bobine (le courant n'est alors plus proportionnel à l'écart de tension)
ou par paliers de tension --- chaque palier de $E/n$ ne perd que
$\tfrac12 C(E/n)^2$.
\end{solution}

\begin{solution}{exo:b1:transient-regimes:7}
$\delta = \ln(2.0/1.5) = 0.29$ (et $\ln(1.5/1.12) = 0.29$~: cohérent)~;
$Q \approx \pi/\delta = 11$. $T = \qty{0.40}{ms}$, donc $\omega_0 \approx
2\pi/T = \qty{1.57e4}{rad/s}$ (la correction $1/8Q^2$ vaut $0.1\%$).
$L = 1/(\omega_0^2 C) = \qty{4.1}{mH}$~; $R = \omega_0 L/Q = \qty{5.8}{\ohm}$.
\end{solution}

\begin{solution}{exo:b1:transient-regimes:8}
$u_C \approx U_0\eu^{-t/\tau}\cos\omega_0 t$~; $i = -C\,\dd u_C/\dd t
\approx CU_0\omega_0\sin\omega_0 t$ (pour $t \ll \tau$), de crête $U_0\sqrt{C/L}
= 100 \times 0.1 = \qty{10}{A}$, atteinte un quart de période après la
fermeture, quand $u_C = 0$~: toute l'énergie ($\tfrac12 CU_0^2 = \qty{50}{mJ}$)
est alors magnétique, $\tfrac12 Li^2 = \qty{50}{mJ}$.
\end{solution}

\begin{solution}{exo:b1:transient-regimes:9}
$\omega_0 = \sqrt{k/m} = \qty{7.1}{rad/s}$, $f_0 = \qty{1.1}{Hz}$~;
$\alpha_c = 2\sqrt{km} = 2\sqrt{8\times10^6} = \qty{5.7e3}{N.s/m}$.
Pour $\alpha = 2000$~: $Q = \sqrt{km}/\alpha = 2828/2000 = 1.4$, \omterm{thm:b1:transient-regimes:regimes}{régime pseudo-périodique}~;
$T = 2\pi/(\omega_0\sqrt{1 - 1/4Q^2}) = 0.89/0.935 = \qty{0.95}{s}$~;
$\tau = 2Q/\omega_0 = \qty{0.40}{s}$.
\end{solution}

\begin{solution}{exo:b1:transient-regimes:10}
$i = I_0\eu^{-t/\tau}$, $\tau = L/R = \qty{5.0}{ms}$~; tension de crête
$RI_0 = \qty{200}{V}$ en $t = 0^+$~; énergie $\tfrac12 LI_0^2 = \qty{1.0}{J}$
dissipée dans $R$. Sans la résistance, le courant n'a plus de chemin~:
$\dd i/\dd t$ devient énorme, la tension de la bobine atteint des
kilovolts, et un arc jaillit à l'interrupteur.
\end{solution}

\begin{solution}{exo:b1:transient-regimes:11}
Vu de $C$~: $E_{\mathrm{Th}} = ER_2/(R_1 + R_2) = \qty{8.0}{V}$,
$R_{\mathrm{Th}} = R_1 \parallel R_2 = \qty{6.67}{k\ohm}$~; $\tau =
R_{\mathrm{Th}}C = \qty{6.7}{ms}$~; $u_C = 8.0(1 - \eu^{-t/\tau})$ V.
\end{solution}

\begin{solution}{exo:b1:transient-regimes:12}
Racine double $-\omega_0$~: $x = x_\infty + (A + Bt)\eu^{-\omega_0 t}$~;
$x(0) = 0$ donne $A = -x_\infty$, $x'(0) = 0$ donne $B = \omega_0 A$~:
$x = x_\infty[1 - (1 + \omega_0 t)\eu^{-\omega_0 t}]$. Pour $95\%$~:
$(1 + s)\eu^{-s} = 0.05$, $s \approx 4.7$~: $t_{95} = 4.7/\omega_0$. Pour
$Q = 5$, l'enveloppe $\eu^{-t/\tau}$ avec $\tau = 10/\omega_0$ demande
$3\tau = 30/\omega_0$~; pour $Q = 0.1$, la racine lente vaut $\abs{r_+} \approx
\omega_0 Q$ et $3/\abs{r_+} = 30/\omega_0$. Les deux sont six fois plus
lents~: l'amortissement critique est le retour sans dépassement le plus
rapide, exactement ce que veulent un appareil de mesure, un ferme-porte
ou une suspension.
\end{solution}

\begin{solution}{pb:b1:transient-regimes:1}
\textbf{1.} $E = L\,\dd i/\dd t + Ri + u_C$, $i = C\,\dd u_C/\dd t$~:
$LC\,u_C'' + RC\,u_C' + u_C = E$.

\textbf{2.} $u_C'' + (\omega_0/Q)u_C' + \omega_0^2u_C = \omega_0^2E$ avec
$\omega_0 = 1/\sqrt{LC}$, $Q = \sqrt{L/C}/R$.

\textbf{3.} $\omega_0 = \qty{3.16e4}{rad/s}$, $f_0 = \qty{5.03}{kHz}$,
$R_c = 2\sqrt{L/C} = \qty{632}{\ohm}$.

\textbf{4.} $r^2 + (\omega_0/Q)r + \omega_0^2 = 0$, $\Delta = \omega_0^2(1/Q^2
- 4)$~: $Q < \tfrac12$ apériodique, $Q = \tfrac12$ critique, $Q > \tfrac12$
pseudo-périodique.

\textbf{5.} $\tau = 2Q/\omega_0$, $\omega = \omega_0\sqrt{1 - 1/4Q^2}$,
$u_C = E + \eu^{-t/\tau}(A\cos\omega t + B\sin\omega t)$.

\textbf{6.} $Q = 3.16$, $\tau = \qty{0.20}{ms}$, $T = 2\pi/\omega =
\qty{0.20}{ms}$, environ trois oscillations.

\textbf{7.} $u_C(0^+) = U_0$ ($u_C$ continue~: courant fini)~;
$u_C'(0^+) = i(0^+)/C = 0$ ($i_L$ continue et nulle avant).

\textbf{8.} $m\ddot x = -k(x + x_{\text{eq}}) - \alpha\dot x + mg$ avec
$kx_{\text{eq}} = mg$~: $m\ddot x + \alpha\dot x + kx = 0$.

\textbf{9.} $\omega_0 = \sqrt{k/m}$, $Q = \sqrt{km}/\alpha$~; $x \leftrightarrow
q$, $\dot x \leftrightarrow i$, $m \leftrightarrow L$, $\alpha \leftrightarrow
R$, $k \leftrightarrow 1/C$.

\textbf{10.} $\omega_0 = \sqrt{22000/350} = \qty{7.93}{rad/s}$, $f_0 =
\qty{1.26}{Hz}$~; $\alpha_c = 2\sqrt{km} = \qty{5.55e3}{N.s/m}$.

\textbf{11.} $Q = 2775/4000 = 0.69$, $\xi = 0.72$~: pseudo-périodique,
juste en deçà du critique.

\textbf{12.} $\tau = 2Q/\omega_0 = \qty{0.18}{s}$~; $\omega = \omega_0\sqrt{1
- \xi^2} = \qty{5.5}{rad/s}$, $T = \qty{1.1}{s}$.

\textbf{13.} Dépassement $\exp(-\pi \times 0.72/0.69) = \exp(-3.3) = 3.8\%$~:
une petite remontée au-delà de l'équilibre, et c'est fini --- confortable.

\textbf{14.} $R = \sqrt{L/C}/Q = 316/0.69 = \qty{455}{\ohm}$.

\textbf{15.} $Q = 2775/1200 = 2.3$, $\xi = 0.22$~; $\omega = 0.976\omega_0$,
$T = \qty{0.81}{s}$~; $\tau = 2Q/\omega_0 = \qty{0.58}{s}$.

\textbf{16.} $\eu^{-T/\tau} = \eu^{-1.39} = 0.25$~: chaque rebond vaut le
quart du précédent~; $0.25^3 = 1.6\%$~: trois rebonds visibles.

\textbf{17.} $x(0) = -x_0$, $\dot x(0) = 0$~; énergie $\tfrac12 kx_0^2 =
0.5 \times 22000 \times 0.0025 = \qty{28}{J}$, dissipée en chaleur dans
l'huile de l'amortisseur.

\textbf{18.} $v \approx x_0\omega_0 = 0.05 \times 7.93 = \qty{0.40}{m/s}$~;
force dans l'amortisseur $\alpha v \approx 1200 \times 0.40 = \qty{480}{N}$.

\textbf{19.} À $Q = \tfrac12$, le retour est en $(1 + \omega_0 t)\eu^{-\omega_0
t}$, qui rampe asymptotiquement~: les derniers millimètres sont longs.
Un système légèrement sous-amorti ($\xi \approx 0.7$) dépasse de
quelques pour cent mais entre plus tôt, et reste, dans une bande de tolérance.

\textbf{20.} $\exp(-\pi \times 0.22/0.976) = \exp(-0.70) = 50\%$, contre
$3.8\%$~: la voiture usée rebondit à mi-hauteur au-delà de l'équilibre.

\textbf{21.} $\delta = \ln 4 = 1.39$ (deux fois, cohérent)~; $\xi =
\delta/\sqrt{4\pi^2 + \delta^2} = 1.39/6.43 = 0.215$, $Q = 1/2\xi = 2.3$~:
la valeur de l'amortisseur usé.

\textbf{22.} $-\ln 0.038 = 3.27 = \pi\xi/\sqrt{1 - \xi^2}$, donc $\xi =
3.27/\sqrt{\pi^2 + 3.27^2} = 0.72$, $Q = 0.69$.

\textbf{23.} L'énergie est $\propto$ \omterm{def:b1:wave-propagation:signal}{amplitude}$^2$~: fraction restante par
période $0.25^2 = 6\%$, donc $94\%$ dissipés par pseudo-période.

\textbf{24.} $R = 316/2.3 = \qty{137}{\ohm}$. Dans la variable $\omega_0 t$,
l'équation canonique ne contient plus que $Q$~: deux systèmes de même
$Q$ tracent la même courbe.

\textbf{25.} Le \omterm{def:b1:transient-regimes:secondorder}{facteur de qualité} $Q$ (ou $\xi$) fixe à lui seul la
forme~; un concepteur vise $Q \approx 0.7$ ($\xi \approx 0.7$)~; plus
d'un rebond visible signifie un $Q$ nettement supérieur à $1$~:
l'amortisseur est usé.
\end{solution}
